%% kapa: exam=1 round=35 year=2024 period=1 form=A subject=civil_law subject_ko=민법 seq=22 number=22
%% source: Q-Net 감정평가사 기출문제 (cst00309) artlSeq=5213229 file=raw/1008/5213229/1교시 문제지 원본.pdf
%% conversion: kordoc 4.24.51 -> tools/gen1.py (auto); review=proofread-sonnet; flags=none
\begin{question}{22}
甲이 乙소유 X토지에 권원없이 Y건물을 신축하여 소유하고 있다. 이에 관한 설명으로 옳은 것은? (다툼이 있으면 판례에 따름)
\begin{choices}
  \item 乙은 Y를 관리하는 甲의 직원 A에게 X의 반환청구를 할 수 있다.
  \item 甲이 법인인 경우 乙은 甲의 대표이사 B 개인에게 X의 반환청구를 할 수 있다.
  \item 乙이 甲에게 X의 반환청구를 하여 승소한 경우, 乙은 甲에게 Y에서 퇴거할 것을 청구할 수 있다.
  \item 미등기인 Y를 丙이 매수하여 인도받았다면 乙은 丙을 상대로 건물철거 청구를 할 수 있다.
  \item 乙은 甲에 대한 X의 반환청구권을 유보하고 X의 소유권을 丁에게 양도할 수 있다.
\end{choices}
\end{question}
\answer{4}
